\documentclass[journal]{IEEEtran}
\usepackage{graphicx}
\usepackage{amsmath}
\usepackage{listings}
\usepackage{xcolor}
\usepackage{hyperref}
\usepackage{booktabs}

\lstset{
  language=Python,
  basicstyle=\ttfamily\small,
  keywordstyle=\color{blue},
  commentstyle=\color{gray},
  stringstyle=\color{orange},
  breaklines=true,
  frame=single,
  numbers=left,
  numberstyle=\tiny\color{gray}
}

\title{Noise Filtering for Line Drawing Preservation:\\Iron Man Image}
\author{Neal Kauntia\\ARK Fresher Selection Task 2.3}

\begin{document}
\maketitle

\begin{abstract}
This document describes the approach taken to remove salt noise from a heavily corrupted 
line drawing of Iron Man (iron\_man\_noisy.jpg) while preserving the fine structural edges 
that constitute the figure. The image suffers from high-frequency salt noise — random white 
dots scattered uniformly across the frame — superimposed on a black-background line drawing. 
A pipeline combining grayscale conversion, median filtering, and morphological opening was 
developed. Extensive kernel size tuning was performed at each stage to strike the optimal 
balance between noise removal and edge preservation. The final output (iron\_man\_final.png) 
successfully removes the majority of salt noise while retaining sharp, continuous edge lines.
\end{abstract}

\section{Introduction}
Salt noise in binary-style line drawings presents a unique denoising challenge. Unlike 
natural images where some blurring is tolerable, a line drawing consists entirely of 
thin, high-frequency edges — the very frequency band that most standard smoothing 
filters attack. Any over-smoothing directly degrades the primary signal.

The task is a balancing act: aggressive enough filtering to eliminate the noise, 
yet gentle enough to leave the line strokes intact. This document details the problem, 
all intermediate approaches explored, and the final chosen pipeline with justification 
for each parameter choice.

\section{Problem Statement}
The input image \texttt{iron\_man\_noisy.jpg} is a grayscale line drawing of Iron Man 
corrupted by dense salt noise. The characteristics of the noise and signal are:

\begin{itemize}
    \item \textbf{Noise type:} Salt noise — isolated single-pixel white dots with 
    intensity 255 scattered at high density across the black background.
    \item \textbf{Signal type:} White line strokes (intensity 255) of width 
    approximately 1--3 pixels on a black (0) background.
    \item \textbf{Core challenge:} Both noise and signal occupy the same intensity 
    value (white), so thresholding alone cannot separate them. Spatial context 
    must be exploited — noise pixels are isolated whereas signal pixels form 
    connected curves.
    \item \textbf{Success criteria:} Maximise Signal-to-Noise Ratio (SNR) while 
    minimising structural damage to the line drawing features.
\end{itemize}

\section{Related Work}
Common approaches for salt noise removal include:

\textbf{Gaussian Blur} averages a pixel with its neighbourhood 
using a Gaussian-weighted kernel. While effective for Gaussian noise, it blurs edges 
proportionally to kernel size, making it poorly suited for line drawing preservation.

\textbf{Median Filter} replaces each pixel with the median value 
of its neighbourhood. It is the classical and most effective method for salt-and-pepper 
noise because isolated noise pixels are outvoted by background neighbours, yet connected 
edges survive if the kernel is not too large relative to line width.

\textbf{Morphological Opening} (erosion followed by dilation) 
removes small isolated bright regions (noise) that are smaller than the structuring 
element, while preserving larger connected structures such as lines.

\section{Initial Attempts}

\subsection{Attempt 1: Gaussian Blur Alone}
The first attempt applied Gaussian blur directly to the grayscale image.

\begin{lstlisting}
gray = cv.cvtColor(img, cv.COLOR_BGR2GRAY)
blurred = cv.GaussianBlur(gray, (3, 3), 0.5)
\end{lstlisting}

\textbf{Kernel size tuning:} A $(3\times3)$ kernel with $\sigma=0.5$ reduced noise 
marginally but left many salt pixels intact. Increasing to $(5\times5)$ with $\sigma=1.0$ 
removed more noise but visibly thickened and blurred fine lines. A $(7\times7)$ kernel 
with $\sigma=1.5$ caused unacceptable edge smearing, merging adjacent strokes into blobs. 
Gaussian blur proved fundamentally ill-suited: to kill isolated white dots it must 
average them, which inevitably blurs the white line strokes by the same amount.

\subsection{Attempt 2: Median Blur with Various Kernel Sizes}
Median blur was applied directly to the grayscale image.

\begin{lstlisting}
gray = cv.cvtColor(img, cv.COLOR_BGR2GRAY)
result = cv.medianBlur(gray, ksize)
\end{lstlisting}

\textbf{Kernel size tuning} (kernel must be odd):

\begin{table}[h]
\centering
\caption{Median Blur Kernel Size Comparison}
\begin{tabular}{@{}cll@{}}
\toprule
Kernel Size & Noise Removed & Edge Damage \\
\midrule
$3\times3$ & Partial — sparse dots remain & Minimal \\
$5\times5$ & Good — most dots gone & Slight line thinning \\
$7\times7$ & Excellent noise removal & Noticeable stroke gaps \\
$9\times9$ & Near-complete & Severe — thin lines break \\
\bottomrule
\end{tabular}
\end{table}

A $3\times3$ kernel was insufficient because the noise density was high enough that 
some noise pixels had noise neighbours, so the median was still white. A $5\times5$ 
kernel gave a good trade-off but left residual isolated noise clusters. This observation 
motivated adding a morphological stage.

\subsection{Attempt 3: Morphological Opening Alone}
Opening with various kernel sizes was applied to the median-filtered image.

\begin{lstlisting}
kernel = cv.getStructuringElement(cv.MORPH_RECT, (k, k))
opened = cv.morphologyEx(gray, cv.MORPH_OPEN, kernel)
\end{lstlisting}

\textbf{Kernel size tuning:}

\begin{table}[h]
\centering
\caption{Morphological Opening Kernel Size Comparison}
\begin{tabular}{@{}cll@{}}
\toprule
Kernel Size & Effect on Noise & Effect on Lines \\
\midrule
$2\times2$ & Removes truly isolated pixels & Minimal damage \\
$3\times3$ & Removes small clusters & Thin lines slightly eroded \\
$4\times4$ & Aggressive noise removal & Lines begin to fragment \\
$5\times5$ & Very aggressive & Fine strokes destroyed \\
\bottomrule
\end{tabular}
\end{table}

A $2\times2$ kernel was found optimal for the morphological stage: it is large enough 
to erode isolated single-pixel salt noise (which cannot survive a $2\times2$ erosion) 
yet small enough to leave connected line strokes intact, since lines of width $\geq2$ 
pixels survive erosion by a $2\times2$ element.

\section{Final Approach}

The final pipeline is a three-stage process:

\begin{enumerate}
    \item \textbf{Grayscale conversion} — reduces the problem to a single channel, 
    consistent with the binary nature of the line drawing.
    \item \textbf{Median blur with $3\times3$ kernel} — removes the bulk of isolated 
    salt pixels while minimally affecting connected line pixels.
    \item \textbf{Morphological opening with $2\times2$ rectangular kernel} — 
    cleans residual isolated or paired noise pixels that survived the median filter.
\end{enumerate}

\begin{lstlisting}
import cv2 as cv

img = cv.imread("iron_man_noisy.jpg")
gray = cv.cvtColor(img, cv.COLOR_BGR2GRAY)

# Stage 1: Median filter - kills isolated salt pixels
medfil = cv.medianBlur(gray, 3)

# Stage 2: Morphological opening - removes residual noise clusters
kernel = cv.getStructuringElement(cv.MORPH_RECT, (2, 2))
morphopen = cv.morphologyEx(medfil, cv.MORPH_OPEN, kernel)

cv.imwrite("iron_man_final.png", morphopen)
\end{lstlisting}

\textbf{Why median kernel $3\times3$:} The noise consists of isolated single pixels. 
In a $3\times3$ neighbourhood of a noise pixel surrounded by black background, 
the median of 9 values (8 black + 1 white) is black — the noise pixel is eliminated. 
A $5\times5$ kernel was not needed and introduced more edge damage than necessary.

\textbf{Why opening kernel $2\times2$:} After median filtering, remaining noise 
tends to be in pairs or small clusters that survived because two adjacent noise pixels 
can outlast a $3\times3$ median. A $2\times2$ opening erodes these clusters. 
MORPH\_RECT was chosen over MORPH\_ELLIPSE because at size $2\times2$ both are 
equivalent, and RECT has no approximation rounding at small sizes.

\textbf{Why grayscale first:} The original image is a black-and-white line drawing 
stored as BGR. Converting to grayscale reduces memory and processing to one channel 
without any loss of information relevant to the denoising task.

\section{Results and Observations}
The final output image shows a clean line drawing with the vast majority of salt 
noise removed. Key observations:

\begin{itemize}
    \item Fine detail lines (e.g., the face visor, suit panel borders) are preserved 
    with minimal thinning or breakage.
    \item The black background is fully restored — no residual white speckle visible 
    at normal viewing scale.
    \item A small number of noise clusters in very high-density noise regions may 
    persist, but the perceptual quality of the figure is fully restored.
    \item Compared to Gaussian blur, which produced a hazy, grey-gradient result 
    even at small kernel sizes, the median + opening pipeline preserves the binary 
    character of the image — lines remain sharp white, background remains black.
\end{itemize}

The choice of median filtering over Gaussian is justified by its fundamental property: 
for a pixel to change value under median filtering, more than half its neighbourhood 
must differ from it. An isolated noise pixel always fails this test; a connected line 
pixel always passes it (as long as the kernel is not larger than the line width).

\section{Future Work}
\begin{itemize}
    \item \textbf{Adaptive median filtering:} Could handle varying noise density 
    across the image by dynamically increasing the kernel size only in high-noise regions.
    \item \textbf{Line Drawing Learning Model} A classical machine learning model — without any deep learning — could be trained on a dataset of line drawings with known noise levels to automatically predict the optimal median blur kernel size for a new unseen drawing.
    
\end{itemize}

\section{Conclusion}
A lightweight two-stage pipeline of median blur ($3\times3$) followed by 
morphological opening ($2\times2$) was found to be the optimal approach for 
removing salt noise from the Iron Man line drawing. The key insight is that 
salt noise is spatially isolated — a property that median filtering exploits 
directly — while Gaussian blurring is inappropriate because it treats signal 
and noise as equivalent frequency content. The final output successfully 
restores the clean line drawing with minimal structural damage.

\begin{thebibliography}{1}

\bibitem{opencv_freecodecamp}
Beau Carnes, https://www.freecodecamp.org/news/opencv-full-course/

\bibitem{opencv_docs}
https://docs.opencv.org/

\bibitem{geeksforgeeks}
https://www.geeksforgeeks.org/

\bibitem{GitHubCommands}
Adam Lusted (etoxin), https://gist.github.com/etoxin/1acb257550b1de60fe99
\end{thebibliography}

\end{document}